\section{Model Uncertainty \& Ensembles}
% Pays off the slide-13 "can reason about model uncertainty" promise and
% makes the PETS / Nagabandi legend on the case-study slide legible.

\begin{frame}{Why Model Uncertainty Matters}
\begin{itemize}
    \item Recall the cliff: the model was \emph{confidently wrong} off the data
\end{itemize}
\end{frame}

\begin{frame}{Why Model Uncertainty Matters}
\definecolor{mbgreen}{RGB}{122,173,75}
\definecolor{mborange}{RGB}{237,148,72}
\definecolor{mbblue}{RGB}{79,129,189}
\begin{itemize}
    \item Recall the cliff: the model was \emph{confidently wrong} off the data
    \item A single model gives one point prediction, no signal that it is guessing
\end{itemize}
\vspace{0.3em}
\centering
\begin{tikzpicture}[font=\scriptsize, >={Stealth[length=2.2mm]}]
    \fill[mbblue!12] (0,-0.05) rectangle (2.6,2.0);
    \node[mbblue!60!black] at (1.3,1.85) {seen data};
    \foreach \x/\y in {0.4/0.7, 0.8/0.95, 1.2/1.1, 1.6/1.0, 2.0/1.2, 2.4/1.05}
        \fill[mbblue] (\x,\y) circle (1.3pt);
    % predictions agree on-data, fan out off-data
    \draw[mbgreen, thick] (2.6,1.1) .. controls (3.6,1.2) .. (5.4,1.6);
    \draw[mborange, thick] (2.6,1.1) .. controls (3.6,1.0) .. (5.4,0.3);
    \draw[mbblue, thick] (2.6,1.1) .. controls (3.6,1.1) .. (5.4,2.0);
    \draw[densely dashed] (2.6,-0.05) -- (2.6,2.0);
    \node[align=center] at (4.2,-0.45) {off-data: the models \emph{disagree}};
\end{tikzpicture}
\end{frame}

\begin{frame}{Ensembles: Disagreement as Uncertainty}
\begin{itemize}
    \item Train an \textbf{ensemble} of $N$ models (different initializations,
          bootstrapped data)
\end{itemize}
\end{frame}

\begin{frame}{Ensembles: Disagreement as Uncertainty}
\begin{itemize}
    \item Train an \textbf{ensemble} of $N$ models (different initializations,
          bootstrapped data)
    \item On data they \emph{agree}; off data they \emph{disagree}
    \item That disagreement measures \textbf{epistemic uncertainty} $u(s,a)$: a
          per state--action score for uncertainty due to \emph{not having seen
          enough data}. It is high off-data, and it \emph{shrinks} as you collect
          more (unlike the environment's inherent noise, which never does)
    \item Plan conservatively where $u(s,a)$ is high: don't trust a plan that
          exploits one model's fantasy
\end{itemize}
\end{frame}

\begin{frame}{PETS: Planning with an Ensemble}
\begin{itemize}
    \item Put the two pieces together:
    \begin{itemize}
        \item \textbf{Planner}: CEM (from the previous section)
        \item \textbf{Model}: an ensemble of probabilistic dynamics models
    \end{itemize}
\end{itemize}
\end{frame}

\begin{frame}{PETS: Planning with an Ensemble}
\begin{itemize}
    \item Put the two pieces together:
    \begin{itemize}
        \item \textbf{Planner}: CEM (from the previous section)
        \item \textbf{Model}: an ensemble of probabilistic dynamics models
    \end{itemize}
    \item Evaluate each candidate plan by propagating it through \emph{sampled}
          models and averaging the return
    \item \textbf{PETS} (Chua et al., 2018): this is the ``CEM-based planner''
          in the case study; Nagabandi's method is the same idea \emph{without}
          the ensemble
\end{itemize}
\end{frame}
